\documentclass[10pt,a4paper]{amsart}
\usepackage{amsmath,amssymb,amsthm,graphicx}
\usepackage{verbatim}
\newenvironment{reference}{\comment}{\endcomment}
%\newenvironment{reference}{}{}
\newenvironment{slogan}{\comment}{\endcomment}
\newenvironment{history}{\comment}{\endcomment}

% For commutative diagrams we use Xy-pic
\usepackage[all]{xy}

% We use 2cell for 2-commutative diagrams.
\xyoption{2cell}
% Theorem environments.
%
\theoremstyle{plain}
\newtheorem{theorem}[subsection]{Theorem}
\newtheorem{proposition}[subsection]{Proposition}
\newtheorem{lemma}[subsection]{Lemma}

\theoremstyle{definition}
\newtheorem{definition}[subsection]{Definition}
\newtheorem{example}[subsection]{Example}
\newtheorem{exercise}[subsection]{Exercise}
\newtheorem{situation}[subsection]{Situation}

\theoremstyle{remark}
\newtheorem{remark}[subsection]{Remark}
\newtheorem{remarks}[subsection]{Remarks}

\numberwithin{equation}{subsection}

% Macros
%
\def\lim{\mathop{\mathrm{lim}}\nolimits}
\def\colim{\mathop{\mathrm{colim}}\nolimits}
\def\Spec{\mathop{\mathrm{Spec}}}
\def\Hom{\mathop{\mathrm{Hom}}\nolimits}
\def\Ext{\mathop{\mathrm{Ext}}\nolimits}
\def\SheafHom{\mathop{\mathcal{H}\!\mathit{om}}\nolimits}
\def\SheafExt{\mathop{\mathcal{E}\!\mathit{xt}}\nolimits}
\def\Sch{\mathit{Sch}}
\def\Mor{\mathop{\mathrm{Mor}}\nolimits}
\def\Ob{\mathop{\mathrm{Ob}}\nolimits}
\def\Sh{\mathop{\mathit{Sh}}\nolimits}
\def\NL{\mathop{N\!L}\nolimits}
\def\CH{\mathop{\mathrm{CH}}\nolimits}
\def\proetale{{pro\text{-}\acute{e}tale}}
\def\etale{{\acute{e}tale}}
\def\QCoh{\mathit{QCoh}}
\def\Ker{\mathop{\mathrm{Ker}}}
\def\Im{\mathop{\mathrm{Im}}}
\def\Coker{\mathop{\mathrm{Coker}}}
\def\Coim{\mathop{\mathrm{Coim}}}

% Boxtimes
%
\DeclareMathSymbol{\boxtimes}{\mathbin}{AMSa}{"02}

%
% Macros for moduli stacks/spaces
%
\def\QCohstack{\mathcal{QC}\!\mathit{oh}}
\def\Cohstack{\mathcal{C}\!\mathit{oh}}
\def\Spacesstack{\mathcal{S}\!\mathit{paces}}
\def\Quotfunctor{\mathrm{Quot}}
\def\Hilbfunctor{\mathrm{Hilb}}
\def\Curvesstack{\mathcal{C}\!\mathit{urves}}
\def\Polarizedstack{\mathcal{P}\!\mathit{olarized}}
\def\Complexesstack{\mathcal{C}\!\mathit{omplexes}}
% \Pic is the operator that assigns to X its picard group, usage \Pic(X)
% \Picardstack_{X/B} denotes the Picard stack of X over B
% \Picardfunctor_{X/B} denotes the Picard functor of X over B
\def\Pic{\mathop{\mathrm{Pic}}\nolimits}
\def\Picardstack{\mathcal{P}\!\mathit{ic}}
\def\Picardfunctor{\mathrm{Pic}}
\def\Deformationcategory{\mathcal{D}\!\mathit{ef}}
\usepackage{cleveref}
\crefname{originallemma}{original lemma}{original lemmas}
\Crefname{originallemma}{Original lemma}{Original lemmas}
\makeatletter
\newlabel{algebra-lemma-graded-silly}{{00JS}{source}{}{sourcepointer-1}{}}
\newlabel{algebra-lemma-graded-silly@cref}{{[originallemma][0][]00JS}{[1][0][]source}{}{}{}}
\makeatother
\crefname{originalreference}{original reference}{original references}
\Crefname{originalreference}{Original reference}{Original references}
\makeatletter
\newlabel{definition-quasi-affine}{{01P6}{source}{}{sourcepointer-2}{}}
\newlabel{definition-quasi-affine@cref}{{[originalreference][0][]01P6}{[1][0][]source}{}{}{}}
\makeatother
\crefname{originalreference}{original reference}{original references}
\Crefname{originalreference}{Original reference}{Original references}
\makeatletter
\newlabel{proposition-characterize-ample}{{01Q3}{source}{}{sourcepointer-3}{}}
\newlabel{proposition-characterize-ample@cref}{{[originalreference][0][]01Q3}{[1][0][]source}{}{}{}}
\makeatother
\title{Finitely many points of a quasi-affine scheme or a scheme with an ample invertible sheaf lie in an affine open}
\author{The Stacks Project Authors}
\date{stacks-project:properties.tex:local-snapshot-20261012; GNU FDL 1.2 or later}
\begin{document}
\maketitle
\noindent\textit{Editorial note.} All delivered mathematical text is the author's own, copied from the pinned Stacks Project chapter file properties.tex; no statement or proof step was rewritten, reordered or removed.\par\medskip
\noindent\textit{Editorial note.} The counted proof is the author's complete argument: the reduction of the four cases to the projective case, the covering of the finite set by finitely many standard opens avoiding the complement, and the graded argument merging them into a single standard open whose intersection with X is affine.\par\medskip
\noindent\textit{Editorial note.} The label lemma-ample-finite-set-in-affine names the counted statement, and the proof environment immediately following it in the source is delivered verbatim as its proof.\par\medskip
\noindent\textit{Editorial note.} The counted statement is carried by the author's own lemma environment in the chapter source.\par\medskip
\noindent\textit{Editorial note.} Pre-open gate: python operations/agent-authoring/tools/precard\_gate.py --source properties.tex --label lemma-ample-finite-set-in-affine --start 5460 --end 5507 --local-id 28036 returned RESULT=OK (L28036 free, occupancy FREE, public interval clear). Candidate taken from sources/stacks-properties/candidates-halfA.json.\par\medskip
\noindent\textit{Editorial note.} The proof contains the sentence 'We urge the reader to prove case (2) directly for themselves.' — a reading suggestion for one of the four hypotheses, while the proof itself reduces that case and the others to the projective case, so the statement is counted and the sentence is recorded here verbatim.\par\medskip
\noindent\textit{Editorial note.} The statement and proof are copied from the pinned chapter file; the cited inputs are kept as external pointers with their canonical Stacks tags.\par\medskip
\noindent\textit{Allowed standard prerequisites (editorial).}\begin{itemize}
\item Description of quasi-affine schemes and of schemes with an ample invertible sheaf. Same chapter, the definition of a quasi-affine scheme, the identification of an affine scheme with a projective spectrum of a polynomial ring, and the proposition characterizing schemes with an ample invertible sheaf as open subschemes of projective spectra. Supplies the reduction of cases (1), (2) and (3) to case (4).
\item Standard opens form a basis of the projective spectrum. Algebra, the description of the standard opens D\_+(f) for homogeneous elements as a basis of the topology of the projective spectrum. Supplies the covering of the finite set by finitely many standard opens avoiding the closed complement.
\item Existence of a homogeneous element avoiding finitely many primes. Algebra, the lemma producing a homogeneous element of a graded ideal which avoids finitely many primes not containing the ideal. Supplies the single homogeneous element f whose standard open contains the finite set and avoids the complement.
\end{itemize}\medskip
\ifdefined\hypertarget\hypertarget{sourcepointer-1}{}\fi
\noindent\textit{Source reference.} Original lemma 00JS: sources/stacks-chow/tags.txt, source lines 726--726. Prerequisite pointer invoked at proof line(s) 5501; it names a standard fact of a sibling chapter that the counted argument uses, and every step of the counted proof built on it is retained verbatim in the delivered body.\par\medskip
\ifdefined\hypertarget\hypertarget{sourcepointer-2}{}\fi
\noindent\textit{Source reference.} Original reference 01P6: sources/stacks-properties/pinned-original-source.tex, source lines 2624--2624. Prerequisite pointer invoked at proof line(s) 5476; it names a standard fact of the same chapter that the counted argument uses, and every step of the counted proof built on it is retained verbatim in the delivered body.\par\medskip
\ifdefined\hypertarget\hypertarget{sourcepointer-3}{}\fi
\noindent\textit{Source reference.} Original reference 01Q3: sources/stacks-properties/pinned-original-source.tex, source lines 4774--4774. Prerequisite pointer invoked at proof line(s) 5480; it names a standard fact of the same chapter that the counted argument uses, and every step of the counted proof built on it is retained verbatim in the delivered body.\par\medskip
\begin{lemma}
\label{lemma-ample-finite-set-in-affine}
Let $X$ be a scheme. Assume either
\begin{enumerate}
\item The scheme $X$ is quasi-affine.
\item The scheme $X$ is isomorphic to a locally closed subscheme
of an affine scheme.
\item There exists an ample invertible sheaf on $X$.
\item The scheme $X$ is isomorphic to a locally closed subscheme
of $\text{Proj}(S)$ for some graded ring $S$.
\end{enumerate}
Then for any finite subset $E \subset X$ there exists an
affine open $U \subset X$ with $E \subset U$.
\end{lemma}

\begin{proof}
By Properties, Definition \ref{definition-quasi-affine}
a quasi-affine scheme is a quasi-compact open subscheme
of an affine scheme. Any affine scheme $\Spec(R)$ is isomorphic to
$\text{Proj}(R[X])$ where $R[X]$ is graded by setting $\deg(X) = 1$.
By Proposition \ref{proposition-characterize-ample}
if $X$ has an ample invertible sheaf then $X$ is isomorphic to an open
subscheme of $\text{Proj}(S)$ for some graded ring $S$.
Hence, it suffices to prove the lemma in case (4).
(We urge the reader to prove case (2) directly for themselves.)

\medskip\noindent
Thus assume $X \subset \text{Proj}(S)$ is a locally closed subscheme
where $S$ is some graded ring. Let $T = \overline{X} \setminus X$.
Recall that the standard opens $D_{+}(f)$ form a basis of the
topology on $\text{Proj}(S)$. Since $E$ is finite we may choose finitely many
homogeneous elements $f_i \in S_{+}$ such that
$$
E \subset
D_{+}(f_1) \cup \ldots \cup D_{+}(f_n) \subset
\text{Proj}(S) \setminus T
$$
Suppose that $E = \{\mathfrak p_1, \ldots, \mathfrak p_m\}$
as a subset of $\text{Proj}(S)$.
Consider the ideal $I = (f_1, \ldots, f_n) \subset S$.
Since $I \not \subset \mathfrak p_j$ for all $j = 1, \ldots, m$
we see from Algebra, Lemma \ref{algebra-lemma-graded-silly} that
there exists a homogeneous element $f \in I$, $f \not \in \mathfrak p_j$
for all $j = 1, \ldots, m$. Then $E \subset D_{+}(f) \subset
D_{+}(f_1) \cup \ldots \cup D_{+}(f_n)$. Since $D_{+}(f)$ does not
meet $T$ we see that $X \cap D_{+}(f)$ is a closed subscheme of the
affine scheme $D_{+}(f)$, hence is an affine open of $X$ as desired.
\end{proof}

\end{document}
